\documentclass[a4paper,10pt,oldfontcommands]{memoir}
\usepackage[headings]{fullpage}
\usepackage[utf8]{inputenc}
\usepackage[T1]{fontenc}
\usepackage{tgpagella}
\usepackage{graphicx}

\DeclareGraphicsExtensions{.pdf,.eps,.png,.jpg}

% formatação do título dos capítulos
\usepackage{titlesec}
\titleformat{\chapter}[hang]
{\normalfont\huge\bfseries}{\chaptertitlename\ \thechapter:}{10pt}{}

\usepackage[brazilian]{babel}
\usepackage[hyphens]{url}
\usepackage{color}
\usepackage{multirow}
% indenta primeira linha de seções
\usepackage{indentfirst}
\usepackage{multicol}
\usepackage{dblfloatfix}
\usepackage{float}
\usepackage{wrapfig}
\usepackage{enumitem}
\setlist[itemize]{leftmargin=*}
\setlist[enumerate]{leftmargin=*}
\setlist[description]{leftmargin=*}
% microtipografia (menos overfull hbox)
\usepackage[final,protrusion=true,expansion]{microtype}
% páginas em branco sem cabeçalho ou rodapé
\usepackage{emptypage}
% figuras lado a lado
\usepackage{subcaption}
% unidades formatadas corretamente
\usepackage{siunitx}
\usepackage[breaklinks]{hyperref}
\hypersetup{
    bookmarks=true,
    pdftitle={Manual d* Bix* \the\year},
    pdfauthor={Centro Acadêmico da Computação - CACo},
    hidelinks
}

% numera somente até o nível de section, mostra no índice até subsection
\setcounter{secnumdepth}{1}
\setcounter{tocdepth}{2}

% indenta subsections no sumário
\usepackage{tocloft}
\cftsetindents{subsection}{5em}{0cm}

% cabeçalho e rodapé
\nouppercaseheads

\makeevenhead{headings}{Manual d* Bix* \the\year}{}{\rightmark}
\makeevenfoot{headings}{\thepage}{}{Centro Acadêmico da Computação}
\makeevenfoot{plain}{\thepage}{}{}

\makeoddhead{headings}{\leftmark}{}{Manual d* Bix* \the\year}
\makeoddfoot{headings}{CACo}{}{\thepage}
\makeoddfoot{plain}{}{}{\thepage}

\renewcommand{\headheight}{13pt}
\makeheadrule{headings}{\textwidth}{\normalrulethickness}
\makefootrule{headings}{\textwidth}{\normalrulethickness}{3pt}

\renewcommand{\chaptermark}[1]{\markboth{#1}{}}
\renewcommand{\sectionmark}[1]{\markright{\thesection\ #1}{}}

% Comando para deixar emails e URLs na formatação e com link correto
\newcommand{\email}[1]{\href{mailto:#1}{\nolinkurl{#1}}}
\renewcommand{\url}[1]{\href{http://#1}{\nolinkurl{#1}}}
\newcommand{\urls}[1]{\href{#1}{\nolinkurl{#1}}}

% Environment que cria listas com menor espaçamento vertical
\newenvironment{compactitemize}{
\begin{itemize}[noitemsep]% ,nolistsep] % Não tirar o espaço vertical, é bonito
}{
\end{itemize}
}
\newenvironment{compactenumerate}{
    \begin{enumerate}[noitemsep]% ,nolistsep]
}{
    \end{enumerate}
}

% Definicao do "shrug" ¯\_(ツ)_/¯
\def\shrug{\normalsize\texttt{\raisebox{0.75em}{\char`\_}\char`\\\char`\_\kern
-0.5ex(\kern-0.25ex\raisebox{0.25ex}{\rotatebox{45}{\raisebox{-.75ex}"\kern
-1.5ex\rotatebox{-90})}}\kern-0.5ex)\kern-0.5ex\char`\_/\raisebox{0.75em}
{\char`\_}}}

\begin{document}

\pagestyle{empty}
\frontmatter
\begin{figure}[H]
  \centering
  \includegraphics[width=.85\textwidth]{img/manual_logo.png}
\end{figure}

\cleardoublepage
\begin{figure}[H]
  \centering
  \includegraphics[width=.85\textwidth]{img/manual_logo.png}
  \caption*{\textsf{\huge Seu guia de sobrevivência na universidade e além}}
\end{figure}
\begin{center}

\vfill % o vfill é pra jogar o logo pro final da página

\begin{figure}[H]
    \centering
    \includegraphics[width=.4\textwidth]{img/caco_logo.pdf}
    \caption*{Centro Acadêmico da Computação

Universidade Estadual de Campinas

\the\year}
\end{figure}



\clearpage
\setlength{\parskip}{1em}
Copyright ~ \copyright ~ 2011-\the\year ~ Centro Acadêmico da Computação

\url{www.caco.ic.unicamp.br}

\begin{figure}[h]
\centering
\begin{subfigure}{0.4\textwidth}
  \centering
  \includegraphics[width=\textwidth]{img/cc_logo.pdf}
\end{subfigure}
\qquad
\begin{subfigure}{0.25\textwidth}
  \centering
  \includegraphics[width=\textwidth]{img/by-sa.pdf}
\end{subfigure}
\end{figure}

Este trabalho está licenciado sob a Licença Attribution-ShareAlike 3.0 Brasil
da Creative Commons. Para ver uma cópia desta licença, visite
\url{creativecommons.org/licenses/by-sa/3.0/br}

Para obter uma cópia do código fonte usado para gerar este livro, visite nosso
repositório em \url{github.com/cacounicamp/Manual-do-Bixo}.

Para obter a versão mais atualizada deste manual, acesse
\url{bit.ly/manual_bixx}.

Capa e verso: Rafael Sartori Martins dos Santos

\setlength{\parskip}{0em}
\vfill

Diagramado em \LaTeX\,e composto em \TeX\,Gyre Pagella

Impresso pela Gráfica Central da Unicamp em couchê \SI{230}{g/m^2} na capa e
offset \SI{75}{g/m^2} no miolo.

\end{center}
\cleardoublepage
\pagestyle{plain}
\tableofcontents

\chapter{Prefácio}
\setlength{\parskip}{1em}

Bem-vind*!

Este é o Manual d* Bix* do CACo, o Centro Acadêmico da Computação. Ele foi
elaborado por veteranas e veteranos (alunas, alunos, ex-alunas e ex-alunos) e
inclui diversas dicas para auxiliar a sua sobrevivência neste primeiro ano na
Unicamp. Você provavelmente recebeu a versão impressa nas primeiras semanas de
aula (senão, passe na salinha do CACo -- no IC-3, ao lado do suporte, como
falaremos pelo manual -- para pegar uma).

Nós trabalhamos muito para que as informações presentes aqui sejam úteis,
corretas e atuais, no entanto, não é possível atualizar o manual todo, refazer
as pesquisas todos os anos, então contamos com a ajuda de você, bixete ou bixo!
Se você encontrar qualquer informação incorreta ou desatualizada, nos avise em
\email{caco@ic.unicamp.br}, mande-nos uma mensagem no Facebook ou o atualize
você mesmo, pois ele está esperando sua colaboração em
\url{github.com/cacounicamp/Manual-do-Bixo}!

\subsection{Sobre o uso de substantivos femininos e *'s}

Com certeza vocês já devem ter notado que o escrito na capa do nosso manual é
um pouco incomum, não é?  Acham que foi erro de digitação? Falta de atenção
nossa? Muito pelo contrário!

Com certeza vocês, ex-vestibulandas e ex-vestibulandos (aleluia!), humanas e
humanos do século XXI, estão antenadas e antenados nas lutas sociais que estão
acontecendo a sua volta, e uma delas é a luta pela igualdade de gêneros. Ela se
baseia na ideia de inclusão de todos os gêneros, de forma que todas as pessoas
tenham os mesmos direitos dentro da nossa sociedade. Dentro dessa luta, entra a
neutralidade de gênero que, aplicada no âmbito linguístico e gramatical, é uma
forma de aplicar a demanda de igualdade de gêneros no nosso modo de escrever e
falar, já que a língua portuguesa não é nem um pouco neutra.

A cada ano na Unicamp temos mais e mais diversidade de gênero e isso é ótimo!
No entanto, para facilitar a leitura dos que tem ou não alguma deficiência
visual, reescrevemos os substantivos que antes utilizavam *'s -- denotando a
terminação que você quiser usar, a que era mais confortável -- para o feminino,
assim podem ser lidos por leitores digitais sem problemas, além de enfatizar o
problema da desigualdade de gênero no âmbito da computação. Quando encontrar
algum substantivo que não siga essa ideia, por favor, nos contate pelo e-mail,
pelo Facebook ou até edite você mesmo usando o repositório no GitHub, como
mencionamos acima.

Vale a pena enfatizar que o seu centro acadêmico não liga para suas
características, está aqui para representar todas e todos estudantes, sem
privilegiar qualquer pessoa. Fazemos o possível para tornar todas as decisões o
mais democráticas possíveis e falaremos mais disso durante o manual.

O CACo, como seu centro acadêmico, veio te trazer esse manual com essa capa
linda e todas essas dicas para que você, (seja) ``aluno'', ``aluna'' ou
``alune'', possa sobreviver as próximas semanas!

Aproveite!

\setlength{\parskip}{0em}
% aqui começa o manual
\mainmatter
\pagestyle{headings}
\twocolumn
\chapter{Olá Mundo!}

%%%%% Bem-vindo
\input{01-ola_mundo/bem_vind0x2a.tex}
\clearpage

% mantemos as mensagens do IC e FEEC com 1 coluna
\onecolumn
%%%%% Mensagem do diretor da FEEC
\input{01-ola_mundo/mensagem_feec.tex}
\clearpage

%%%%% Mensagem do diretor do IC
\input{01-ola_mundo/mensagem_ic.tex}
\clearpage

% formatamos o resto do manual com 2 colunas
\twocolumn
%%%%% Para pensar
\input{01-ola_mundo/pra_pensar.tex}

\chapter{Sobrevivendo em Barão}
%%%%% Lugares para morar
\input{02-barao/lugares_para_morar.tex}
\newpage
%%%%% Comida
\input{02-barao/comida.tex}
\newpage
%%%%% Diversão
\input{02-barao/diversao.tex}
\newpage
%%%%% Transporte
\input{02-barao/transporte.tex}
\newpage
%%%%% Outras necessidades
\input{02-barao/outras_necessidades.tex}

\chapter{Estudando na Unicamp}
%%%%% Disciplinas
\input{03-unicamp/disciplinas.tex}
\newpage
%%%%% SIGLAS malucas
\input{03-unicamp/siglas.tex}
\newpage
%%%%% Lugares para estudar
\input{03-unicamp/lugares_para_estudar.tex}
\newpage
%%%%% Cuidado com CR e Reprovação
\input{03-unicamp/cr_reprovacao.tex}
\newpage
%%%%% Pra que que eu estou estudando isso??
\input{03-unicamp/pra_que_estou_estudando_isso.tex}
\newpage
%%%%% Bem-vindo à academia
\input{03-unicamp/bem_vindo_a_academia.tex}
\newpage
%%%%% Modalidades da engenharia de computação
\input{03-unicamp/modalidades.tex}
\newpage
%%%%% Emails acadêmicos
\input{03-unicamp/email.tex}
\newpage
%%%%% Melhores banheiros
\input{03-unicamp/melhores_banheiros.tex}

\chapter{Além da graduação}
%%%%% Centro Acadêmico da Computação
\input{04-alem_da_graduacao/caco.tex}
\newpage
%%%%% Atlética -- AAACEC
\input{04-alem_da_graduacao/aaacec.tex}
\newpage
%%%%% Bateria Valorosa
\input{04-alem_da_graduacao/valorosa.tex}
\newpage
%%%%% Conpec
\input{04-alem_da_graduacao/conpec.tex}
\newpage
%%%%% Computing 4 All
\input{04-alem_da_graduacao/computing4all.tex}
\newpage
%%%%% Enigma
% Este arquivo .tex será incluído no arquivo .tex principal. Não é preciso
% declarar nenhum cabeçalho

\section{Enigma}

O Enigma é uma entidade formada por alunas e alunos da Unicamp com o propósito
de estudar privacidade, segurança e criptografia, além de outros assuntos
relacionados. Inauguramos nossas atividades com um \textit{Capture the Flag} e
uma palestra na Secomp 2018. Desde então, organizamos encontros, mini cursos,
palestras e desafios para incentivar o desenvolvimento desta área na
universidade.

A área de segurança da informação deixou de ser uma disciplina de nicho para
ser um pré-requisito para todo software que envolve conexão com internet e/ou
dados sensíveis. Todo site e programa pode ser vulnerável, e uma das formas de
aprender a se defender, é aprender o ataque. Por isso, em nossas atividades,
exploramos e discutimos vulnerabilidades, bem como meios de mitigá-las. A
privacidade é um direito fundamental do ser humano e que garante um
funcionamento democrático da sociedade. Esse direito tem sido cada vez mais
negligenciado na era digital, com abusos por parte de corporações e governos.
Criptografia é a técnica utilizada para transmitir e armazenar dados de forma
segura e garantida que somente aqueles com os devidos poderes podem vir a
acessá-los.

Todas as pessoas são bem vindas, independentemente do seu curso e nível
técnico! Somos abertos, sem processo seletivo ou hierarquias e acreditamos que
o conhecimento deve ser livre.

\begin{figure}[H]
  \centering
  \includegraphics[width=.24\textwidth]
  {img/alem_da_graduacao/enigma_logo.jpg}
\end{figure}


\begin{compactitemize}
\item Site: \url{enigma.ic.unicamp.br/}
\item Email: \url{enigmaunicamp@tutanota.com}
\item Grupo do Telegram: \url{https://t.me/enigmaunicamp}
\end{compactitemize}

\newpage
%%%%% Grupos e Entidades da Unicamp
\input{04-alem_da_graduacao/grupos_entidades.tex}
\newpage
%%%% Eventos da Unicamp
\input{04-alem_da_graduacao/eventos.tex}
\newpage
%%%%% Serviços da Unicamp
\input{04-alem_da_graduacao/servicos.tex}

\end{document}
